\section{Reduced-order modelling in \texorpdfstring{\rom{}}{rom\_engine}}\label{sec:rom}

This section states the methods in \rom{}. All of them are established; the contribution here is a uniform array-level implementation and the way the pieces compose. We write $n$ for the full dimension and $r\ll n$ for the reduced one, and
assume $\K,\M$ are symmetric with $\K$ positive definite on the free DOFs. The general setting and its literature are in \cite{benner2015survey,hesthaven2016rb,antoulas2005}.

\subsection{Projection-based reduction}\label{sec:rom:proj}
\paragraph{Bases.}
A reduced basis $\bm V\in\mathbb R^{n\times r}$ gives the approximation $\vu\approx\bm V\bm q$. For snapshots $\bm X=[\vu^{(1)},\dots,\vu^{(m)}]$ the proper orthogonal decomposition (POD) is the
truncated singular value decomposition $\bm X\approx\bm U_r\bm\Sigma_r\bm W_r^{\!\top}$, which is the best rank-$r$ approximation in the Frobenius norm \cite{berkooz1993pod}.
For structural models the Euclidean norm mixes translations and rotations, so \texttt{PodBasis.fit(\dots, M=M)} uses the mass inner product: with the Cholesky factor $\M=\bm L\bm L^{\!\top}$,
\begin{equation}
  \tilde{\bm X}=\bm L^{\!\top}\bm X=\tilde{\bm U}\bm\Sigma\bm W^{\!\top},\qquad \bm V=\bm L^{-\top}\tilde{\bm U}_r,\qquad \bm V^{\!\top}\M\bm V=\bm I_r .
  \label{eq:pod}
\end{equation}
The rank is set by a count or by an energy fraction $\sum_{i\le r}\sigma_i^2/\sum_i\sigma_i^2$. Mode shapes from \texttt{solve\_modal} are also valid bases, and a multi-field variant fits separate fields with a block-diagonal weight.

\paragraph{Galerkin projection.}
Requiring the residual to be orthogonal to the basis, $\bm V^{\!\top}(\K\bm V\bm q-\F)=\bm 0$, gives the reduced system
\begin{equation}
  \K_r\bm q=\F_r,\qquad \K_r=\bm V^{\!\top}\K\bm V,\ \ \M_r=\bm V^{\!\top}\M\bm V,\ \ \F_r=\bm V^{\!\top}\F .
  \label{eq:galerkin}
\end{equation}
For symmetric positive definite $\K$ this is the best approximation in $\operatorname{range}(\bm V)$ in the energy norm, the same variational argument that makes the finite element method a Galerkin method. \texttt{GalerkinROM} also solves the reduced
eigenproblem $\K_r\bm\phi=\omega^2\M_r\bm\phi$ and expands results to full coordinates by $\vu=\bm V\bm q$.

\paragraph{Static completion of truncated modes.}
A truncated modal basis omits the static contribution of the discarded modes. Two closed-form remedies (mode acceleration and modal truncation augmentation) \cite{besselink2013comparison} are provided. The static correction for a load $\F$ is
\begin{equation}
  \bm q_{\mathrm{cor}}=\K^{-1}\F-\bm V\bigl(\bm V^{\!\top}\K\bm V\bigr)^{-1}\bm V^{\!\top}\F,
  \label{eq:modeacc}
\end{equation}
which needs one full static solve; it can be added after the fact or appended to the basis after $\M$-orthogonalisation against $\bm V$. An error is raised if the correction already lies in $\operatorname{range}(\bm V)$.

\subsection{Parametric problems and certified error}\label{sec:rom:param}
\paragraph{Affine decomposition and the offline--online split.}
If $\K(\bm\mu)=\sum_{q=1}^{Q}\theta_q(\bm\mu)\K_q$ with fixed matrices $\K_q$ (for example one per material region, each assembled with unit modulus), the reduced matrices are
$\K_r(\bm\mu)=\sum_q\theta_q(\bm\mu)\,\bm V^{\!\top}\K_q\bm V$; the projections $\bm V^{\!\top}\K_q\bm V$ are computed once, and each query costs $O(Qr^2+r^3)$, independent of $n$. The decomposition is exact when the physics is affine in the parameters, which holds for linear elasticity in Young's modulus.
\Cref{alg:offline} shows the split.

\begin{algorithm}[t]
\caption{Offline--online reduced-basis solve with affine parameters}\label{alg:offline}
\begin{algorithmic}[1]
\Require snapshots $\bm X$, affine terms $\{\K_q\}$, load $\F$
\State \textbf{Offline:} $\bm V\gets$ POD$(\bm X)$ \eqref{eq:pod}; $\K_q^r\gets\bm V^{\!\top}\K_q\bm V$ for all $q$; $\F_r\gets\bm V^{\!\top}\F$
\State \textbf{Online, for each $\bm\mu$:}
\State \quad $\K_r\gets\sum_q\theta_q(\bm\mu)\K_q^r$; \ solve $\K_r\bm q=\F_r$; \ $\vu\gets\bm V\bm q$ (only if the full field is needed)
\end{algorithmic}
\end{algorithm}

\paragraph{Frequency domain.}
Harmonic response has $\bm A(\omega)\bm X=\F$ with $\bm A(\omega)=-\omega^2\M+\mathrm i\omega\C+\K$, which is affine in $\{\M,\C,\K\}$ with $\bm\theta(\omega)=[-\omega^2,\ \mathrm i\omega,\ 1]$. \texttt{FrequencyROM} is this decomposition composed with Galerkin projection, so a sweep of thousands of frequencies solves $r\times r$ systems.
A basis can be trained on POD of response snapshots at chosen frequencies.

\paragraph{Weak greedy selection.}
Instead of a uniform frequency grid, \cref{alg:greedy} picks training frequencies where the current basis is weakest, which concentrates full-order solves near resonances. The indicator is the difference between the current reduced response and that of the previous (smaller) model evaluated at every untried frequency, the hierarchical error indicator of the weak-greedy construction \cite{hesthaven2016rb}.

\begin{algorithm}[t]
\caption{Weak-greedy frequency basis}\label{alg:greedy}
\begin{algorithmic}[1]
\Require candidate frequencies $\Omega$, seed size $n_0$, tolerance $\epsilon$, maximum size $r_{\max}$
\State choose $n_0$ seed frequencies; solve full-order at each; build $\bm V$ and the reduced model $\mathcal R_k$
\For{$k=1,2,\dots$ until $r\ge r_{\max}$}
  \State for each untried $\omega\in\Omega$: $\eta(\omega)\gets\|\mathcal R_k(\omega)-\mathcal R_{k-1}(\omega)\|$ \Comment{reduced queries only}
  \State $\omega^\ast\gets\arg\max\eta$;\ \textbf{stop if} $\eta(\omega^\ast)<\epsilon$
  \State solve full-order at $\omega^\ast$; rebuild $\bm V$ and $\mathcal R_{k+1}$
\EndFor
\end{algorithmic}
\end{algorithm}

\paragraph{A posteriori bound.}
For any reduced solution $\bm x_r$, the residual $\bm\rho=\F-\bm A(\omega)\bm x_r$ gives
$\|\bm x-\bm x_r\|\le\|\bm\rho\|/\sigma_{\min}(\bm A(\omega))$,
so a rigorous error bound needs a lower bound on the smallest singular value. \rom{} provides two. The first uses Weyl's inequality $|\sigma_{\min}(\bm A)-\sigma_{\min}(\bm B)|\le\|\bm A-\bm B\|_2$ around reference frequencies
$\omega_j$, applied term by term to the affine form $\bm A=\sum_q\theta_q\bm A_q$, which gives $\sigma_{\min}(\bm A(\omega))\ge\max\Bigl(0,\max_j\bigl[\sigma_{\min}(\bm A(\omega_j))-\sum_q|\theta_q(\omega)-\theta_q(\omega_j)|\,\|\bm A_q\|_2\bigr]\Bigr)$ with reference points chosen greedily. The second is the successive constraint method (SCM) of Huynh et al.\ \cite{huynh2007scm} through a small linear program, applied to the real symmetric positive semi-definite matrix
$\bm A^{\!\mathsf H}\bm A$ so that complex, non-Hermitian operators are covered. The sharper natural-norm variant for non-coercive operators is not implemented. Away from the reference frequencies the bound can fall to zero, which is valid but uninformative; adding references tightens it and never invalidates it.

\subsection{System-theoretic reduction}\label{sec:rom:sys}
With input matrix $\bm B$ (load pattern) and output matrix $\bm C_{\mathrm o}$ (selected DOFs), the second-order system is cast into first-order form with state $[\bm q;\dot{\bm q}]$, giving $\bm E\dot{\bm x}=\bm A\bm x+\bm Bu$, $y=\bm C_{\mathrm o}\bm x$ in two algebraically equivalent forms.
Four families follow the classification of Besselink et al.\ \cite{besselink2013comparison}.

\paragraph{Krylov moment matching.}
Expanding the transfer function $H(s)=\bm C_{\mathrm o}(s\bm E-\bm A)^{-1}\bm B$ about $s_0$, a basis of the Krylov subspace $\mathcal K_k\bigl((\bm A-s_0\bm E)^{-1}\bm E,(\bm A-s_0\bm E)^{-1}\bm B\bigr)$ built by block Arnoldi gives a reduced model that matches the first moments at $s_0$. The approximation is local: accurate near $s_0$ and degrading away from it. A two-sided (Petrov--Galerkin) option builds a second basis from the dual system.
SOAR (second-order Arnoldi) \cite{bai2005soar} matches the same moments from the second-order pencil $H(s)=\bm C_{\mathrm o}(s^2\M+s\C+\K)^{-1}\bm B$ without forming the first-order system, using the two-term recurrence $\bm\Phi_j=\bm A_1\bm\Phi_{j-1}+\bm A_2\bm\Phi_{j-2}$ with
$\bm A_1=-\bm K_0^{-1}(\C+2s_0\M)$, $\bm A_2=-\bm K_0^{-1}\M$, $\bm K_0=\K+s_0\C+s_0^2\M$.

\paragraph{Balanced truncation.}
Balanced truncation \cite{moore1981} solves the two Lyapunov equations
\begin{equation}
  \bm A\bm P+\bm P\bm A^{\!\top}+\bm B\bm B^{\!\top}=\bm 0,\qquad
  \bm A^{\!\top}\bm Q+\bm Q\bm A+\bm C_{\mathrm o}^{\!\top}\bm C_{\mathrm o}=\bm 0,
  \label{eq:lyap}
\end{equation}
balances them to equal diagonal Gramians whose entries are the Hankel singular values $\sigma_1\ge\sigma_2\ge\cdots$, and keeps the $r$ states with the largest $\sigma_i$. It preserves stability and satisfies the a priori bound
$\|G-G_r\|_{\mathcal H_\infty}\le2\sum_{i>r}\sigma_i$ \cite{enns1984,glover1984}; the code returns this bound with the model. Because it uses the ports, balanced truncation is accurate over the whole frequency axis rather than near one point.
Related variants in the package are singular perturbation approximation \cite{liu1989spa}, which preserves the DC gain, frequency-weighted balanced truncation with low-pass and band-pass weights, and optimal Hankel-norm approximation, whose realisation \cite{glover1984} attains $\|G-G_r\|_{\mathrm H}=\sigma_{r+1}$ exactly but is implemented for single-input single-output systems only.
A passivity diagnostic is also provided: for a collocated force input and velocity output, the energy balance $\dot E=-\dot{\bm q}^{\!\top}\C\dot{\bm q}+u\,y\le u\,y$ holds when $\C$ is positive semi-definite, so passivity can be certified in closed form for the full model and checked numerically for the reduced one.

\paragraph{Data-driven identification.}
\texttt{LoewnerROM} identifies natural frequencies, damping ratios and mode shapes from sampled frequency responses alone, using the Loewner pencil built from two disjoint frequency sets \cite{mayo2007loewner}. It never receives $\K$, $\M$ or $\C$, so it applies to measured data.

\subsection{Nonlinear reduction}\label{sec:rom:nl}
\paragraph{Intrusive nonlinear Galerkin.}
Projecting $\M\ddot{\vu}+\C\dot{\vu}+\fint(\vu)=\fext(t)$ gives
$\M_r\ddot{\bm q}+\C_r\dot{\bm q}+\bm V^{\!\top}\fint(\bm V\bm q)=\bm V^{\!\top}\fext(t)$
with full (not diagonal, not mass-normalised) reduced matrices, as in the POD--Galerkin reduction of a planar nonlinear rod \cite{georgiou2005}. The reduced tangent $\K_r$ is computed once from the full tangent at $\vu=\bm 0$,
and the nonlinear part is $\bm f_{\mathrm{nl}}(\bm q)=\bm V^{\!\top}\fint(\bm V\bm q)-\K_r\bm q$. Integration is by Newton--Newmark, Runge--Kutta or a generic ODE solver. The cost per step is still proportional to the full model because $\fint(\bm V\bm q)$ is evaluated at all elements, which motivates hyper-reduction.

\paragraph{Hyper-reduction.}
Energy-conserving sampling and weighting (ECSW) \cite{farhat2015ecsw} approximates the reduced force by a weighted sum over a small element subset $\mathcal E$,
\begin{equation}
  \bm f_r(\bm q)=\sum_{e=1}^{n_e}\bm V_e^{\!\top}\bm f_e(\bm V_e\bm q)\ \approx\ \sum_{e\in\mathcal E}w_e\,\bm V_e^{\!\top}\bm f_e(\bm V_e\bm q),\quad w_e\ge0 ,
  \label{eq:ecsw}
\end{equation}
where $\bm V_e$ are the rows of $\bm V$ for element $e$. The weights solve a non-negative least-squares problem built from training states (\cref{alg:ecsw}) with the active-set solver of Lawson and Hanson \cite{lawson1995lsq}, stopped when the relative residual is below a tolerance, which keeps $\mathcal E$ small.
Positive weights over real element contributions preserve the symmetry and positive definiteness of the reduced tangent.

\begin{algorithm}[t]
\caption{ECSW weight training}\label{alg:ecsw}
\begin{algorithmic}[1]
\Require basis $\bm V$, training states $\{\vu^{(k)}\}_{k=1}^{m}$, element force callbacks $\bm f_e$, tolerance $\epsilon$
\State for each $k$ and $e$: $\bm G_{ke}\gets\bm V_e^{\!\top}\bm f_e(\vu^{(k)}_e)$ (element force at the snapshot's own element DOFs), stack over $k$ into column $e$ of $\bm G$
\State $\bm b\gets\sum_e\bm G_{:e}$ \Comment{exact reduced force}
\State solve $\min_{\bm w\ge0}\|\bm G\bm w-\bm b\|$ by Lawson--Hanson NNLS, stopping when $\|\bm G\bm w-\bm b\|\le\epsilon\|\bm b\|$
\State $\mathcal E\gets\{e:w_e>0\}$
\end{algorithmic}
\end{algorithm}

Discrete empirical interpolation (DEIM) \cite{chaturantabut2010deim} instead approximates the force vector from its values at $m$ interpolation DOFs, $\bm f\approx\bm U(\bm P^{\!\top}\bm U)^{-1}\bm P^{\!\top}\bm f$, where $\bm U$ is a POD basis of force snapshots and $\bm P$ selects the DOFs;
QDEIM selects them by column-pivoted QR of $\bm U^{\!\top}$ \cite{drmac2016qdeim}. Gappy reconstruction is the same interpolation applied to a field. ECSW needs per-element force callbacks; DEIM does not.

\paragraph{Surrogate (implicit-condensation) ROMs.}
In a hybrid reduced model the linear part is intrusive and only the nonlinear modal force is fit from data. In mass-normalised modal coordinates
\begin{equation}
  \ddot{\bm q}+\C_r\dot{\bm q}+\bm\Lambda\bm q+\bm F_{\mathrm{nl}}(\cdot)=\bm F_{\mathrm{ext}}(t),
  \label{eq:hybrid}
\end{equation}
where $\bm\Lambda=\operatorname{diag}(\omega_i^2)$. Two regressions of $\bm F_{\mathrm{nl}}$ are provided, both following published methods:
(i) the implicit condensation and expansion (ICE) method of Hollkamp and Gordon \cite{hollkamp2008ice}, in which $\bm F_{\mathrm{nl}}$ is a quadratic-plus-cubic polynomial in the modal displacements, fit by least squares to full-order solutions under sampled static loads (reviewed in \cite{mignolet2013review}); and
(ii) the multi-fidelity-surrogate method (MFS-NLROM) of He et al.\ \cite{he2023mfsnlrom}, where $\bm F_{\mathrm{nl}}$ is a radial-basis-function interpolant of the \emph{linear} modal displacement,
\begin{equation}
  \bm F_{\mathrm{nl}}(\bm q_l)=\sum_{j=1}^{N}\bm w_j\,\varphi\bigl(\|\bm q_l-\bm q_{l,j}\|\bigr),\qquad \bm W=(\bm\Psi+\varepsilon\,\tfrac{\operatorname{tr}\bm\Psi}{N}\bm I)^{-1}\bm F_{\mathrm{train}},
  \label{eq:rbf}
\end{equation}
with kernel $\varphi$ (multiquadric by default) and a small ridge $\varepsilon$ for conditioning. Because the argument is the linear displacement, evaluating \eqref{eq:rbf} needs no Newton iteration; the polynomial version, which is a function of the nonlinear coordinates, does for static loads. Training data come from applied-load or enforced-displacement strategies. A small neural network can replace the RBF interpolant through the same protocol.
For dynamics, the reduced equations are integrated by Newmark with the nonlinear force lagged by one step and corrected by fixed-point iteration with four iterations (default), by a sign-deviation predictor following He et al., or by Newton. The fixed-point default was
adopted because one-shot lagging was found to accumulate error and diverge on the flat-beam benchmark. Membrane expansion, which adds the in-plane response that flat-beam bending modes lack, follows the ICE paper.

\paragraph{Nonlinear normal modes.}
For the undamped autonomous system, periodic solutions are written as truncated Fourier series $\bm q(t)=\bm Z_0+\sum_{k=1}^{H}[\bm a_k\cos k\omega t+\bm b_k\sin k\omega t]$, and the harmonic-balance equations
\begin{equation}
  \bm h(\bm Z,\omega)=\bm A(\omega)\bm Z+\bm F_{\mathrm{nl}}^{\mathrm{harm}}(\bm Z)=\bm 0
  \label{eq:hbm}
\end{equation}
are solved with the nonlinear force evaluated by alternating frequency--time (AFT) transforms, as in the NNM literature \cite{kerschen2009nnm}. The backbone, frequency against amplitude of a chosen mode, is traced either by natural-parameter continuation in the amplitude (phase fixed by setting the first sine coefficient of the master mode to zero)
or by pseudo-arclength continuation \cite{keller1977} with a secant predictor, which follows folds in amplitude. Damped, forced steady states are instead studied in the time domain with stroboscopic sampling.

\subsection{Component mode synthesis}\label{sec:rom:cms}
Guyan (static) condensation keeps a set of master DOFs and eliminates the rest statically \cite{guyan1965}; it is exact for static loads on the masters, while the condensed mass is approximate. The Craig--Bampton method \cite{craig1968} keeps the interface DOFs $b$ physical and adds the $k$ lowest fixed-interface modes:
\begin{equation}
  \begin{bmatrix}\vu_b\\ \vu_i\end{bmatrix}=\bm T\begin{bmatrix}\vu_b\\ \bm\eta\end{bmatrix},\qquad
  \bm T=\begin{bmatrix}\bm I&\bm 0\\ -\K_{ii}^{-1}\K_{ib}&\bm\Phi_k\end{bmatrix},\qquad
  \K_{ii}\bm\Phi=\M_{ii}\bm\Phi\bm\Omega^2,
  \label{eq:cb}
\end{equation}
with $\bm\Phi_k^{\!\top}\M_{ii}\bm\Phi_k=\bm I$. The reduced matrices are $\bm T^{\!\top}\K\bm T$ and $\bm T^{\!\top}\M\bm T$, whose modal blocks are diagonal by construction and are set exactly. \texttt{couple} joins several substructures by assigning each interface DOF a global index and summing the reduced matrices over the shared coordinates, so substructures can be reduced separately and recombined without re-reduction.
The code requires $\K_{ii}$ positive definite and raises an error for a substructure whose internal DOFs are not constrained by its interface.

\subsection{Validation utilities and experimental components}
\texttt{validate\_rom} compares any reduced predictor with a full-order predictor on held-out inputs and returns relative errors, timing and a pass/fail against a tolerance; \texttt{convergence\_study} and \texttt{select\_basis\_size} give error against reduced size.
The package also contains, as components under development and not evaluated in this paper, a neural-operator residual, a parameterised latent neural ODE, a regularised-projection encoder that is independent of the mesh, an ensemble wrapper for uncertainty, and a differentiable correction of the reduced residual.
